% 190_T0_Korrekturen_En_ch_NEU.tex — Pure correction register (new edition 7 September 2026)
% Replaces the edition of 23 August 2026, which is frozen as 190-Archive-2.
% Full elaborations: 190-Archive-1 (up to R85), 190-Archive-2 (R86–R114).

\providecommand{\QM}{\textbf{[Q]}}

\section*{Preamble}

This document is the ongoing correction register of the FFGFT corpus.
It records \emph{exclusively} two items per entry:
\textbf{(1)} the correction or clarification in brief form and
\textbf{(2)} which document already contains the implemented correction.

Elaborations, derivations, and status notes belong in the target documents
or in the archive editions (Doc.~190-Archive-1: up to R85; Doc.~190-Archive-2: R86–R114).

\textbf{Principle (append-only):} Source documents are \emph{not} revised.
Every declared correction must be visible in an \emph{active} document of the
corpus --- either in a new successor document containing the corrected version,
or as an addendum at the end of the affected document.
An entry that exists only in the archive but appears in no active document
is not complete. The column ``Implemented in'' of the register table names
the active document carrying the correction; it remains empty while that
is still outstanding.

\textbf{Entry types:}
K~$=$~Correction (content error),
P~$=$~Precision (clarification, not an error),
R~$=$~Register entry (revision of a previously booked statement).

\textbf{Markers in the column ``Implemented in'':}

The column does not only contain document numbers but also qualifying
additions that arise from the content of the correction:

\begin{center}
\renewcommand{\arraystretch}{1.3}
\small
\begin{tabular}{@{}p{4.8cm}p{8.7cm}@{}}
\toprule
Marker & Meaning \\
\midrule
\texttt{Doc.~N} & Correction implemented in Doc.~N; status marker absent or not relevant. \\[2pt]
\texttt{Doc.~N \textbf{[B]}} & Correction implemented in Doc.~N and algebraically proven. \\[2pt]
\texttt{Doc.~N \textbf{[K]}} & Correction implemented in Doc.~N and numerically confirmed. \\[2pt]
\texttt{Doc.~N \textbf{[K/B]}} & Correction partly numerically confirmed, partly algebraically proven. \\[2pt]
\texttt{Doc.~N \textbf{[B/S]}} & Correction algebraically established, but one aspect still open. \\[2pt]
\texttt{Doc.~N, open \textbf{[S]}} & Correction declared in Doc.~N, content still open. \\[2pt]
\texttt{Doc.~N (addendum)} & Correction appears only as an addendum at the end of Doc.~N; main text unchanged (append-only). \\[2pt]
\texttt{respective docs} & The correction is implemented in each of the documents listed in column~2; no single target document, but all listed ones. \\[2pt]
\texttt{Doc.~N and respective docs} & Main implementation in Doc.~N; additionally in each of the further listed documents. \\[2pt]
\texttt{script corrected} & Correction applied directly in the verification script; no \TeX{} document affected. \\[2pt]
\texttt{(empty)} & No active document yet carries the correction. Entry open. \\
\bottomrule
\end{tabular}
\renewcommand{\arraystretch}{1.0}
\end{center}

\section{Register Table}

\noindent Full elaborations: Doc.~190-Archive-1 (K1–K7, P1–P43, R1–R85),
Doc.~190-Archive-2 (R86–R114, K-041).

\begingroup\small
\begin{longtable}{@{}p{1.2cm}p{3.4cm}p{5.8cm}p{3.4cm}@{}}
\toprule
\textbf{No.} & \textbf{Concerns} & \textbf{Subject} & \textbf{Implemented in} \\
\midrule\endhead
\bottomrule\endfoot

K1 & 049 (HTML) & $\xi = \lambda_h^2 v^2/(16\pi^3 m_h^2)$; earlier version counted one 2-sphere too many & Doc.~310 \\
K2 & 116 (Tab.) & $r_\tau = 25/9$ & Docs.~016, 046, 116 (De+En) \\
K3 & 018 (Explorer) & $a_e = 2\pi^2/(f \cdot k)$ & Doc.~018 \\
K4 & 267, 268 & Factor~3 $=$ three observable spatial dimensions ($k_\mathrm{obs}^2$ from $n_{1,2,3}$, fourth direction = time unfolding) & Docs.~267, 268 \\
K5 & 275 (script v3) & Three errors in v3; all fixed in v4 (11~June 2026) & Script v4 \\
K6 & 004, 025, 030, 039, 041, 053, 061, 063, 068, 081 & Redshift is achromatic: $1+z=e^{\xi x}$ independent of $\lambda$ & Doc.~280 \\
K7 & rsa/ (14 HTML pages) & Calculation path, labels and method text corrected; no warning boxes, pages describe themselves & rsa/ directly \\
P1 & 116, 158 & Both correct; $0.001\,\%$ for external use & Docs.~116, 158 (addendum) \\
P2 & several & $\hbar$ follows from $\xi$ via $L_0 = \xi\cdot\ell_P$ & Doc.~A080, A085 \textbf{[B]} \\
P3 & 173--176 & Two parameters: $\xi_0$ (geom.) and $\xi_{\text{Higgs}}$ (alg.) & Docs.~173--176 \\
P4 & 180--182 & $L_0$ = geometric fundamental length & Doc.~180 \\
P5 & several & Mass structure from $\xi$; $\sim 1\,\%$ agreement & Doc.~A100 \textbf{[K]} \\
P6 & 116, 158 & Koide only with numerical values (non-closure, Doc.~193) & Doc.~193 \\
P7 & 005, 008, 019, 086, 189, 202 & Scale structure from recursion $\mathcal{R}$ & Doc.~A050 \textbf{[B]} \\
P8 & 005, 008, 012, 014, 041, 060, 133, 141, 159, 181 & Geometric scaling / fractal correction & Doc.~A040 \textbf{[B]} \\
P9 & 005, 006, 042, 046 & Columns ``scaling class'' + ``SM correspondence''; quarks via $m = r\,\xi^p\,v$ & respective docs~\textbf{[K]} \\
P10 & 009 & $K = R_{pe}/(4/3)^7$: empirical anchor; factor $1.314$ implicitly defined & Doc.~009 \\
P11 & 025, 003, 133 & $(1-275/4\cdot\xi)$, $\Delta=0.0002\,\%$; Docs.~003 and 133 use different models & open \textbf{[S]} \\
P12 & 022, 230 & Two observables: cumulative (022) vs.\ elementary $\xi/(2\pi)$ (230) & Docs.~022, 230 \\
P13 & 013, 025, 061 & $\xi_\text{FFGFT}/\xi_\text{UIFT} \approx 414{,}684$; table in 013 reconstructed & Doc.~013 \\
P14a & 041, 268 & Fractal path lengthening (geometric path stretch); energy loss only artefact & Docs.~041, 268 \\
P15 & 061, 041 & Marked as $\Lambda$CDM reference quantity; static universe has no cosmological $z$ & Doc.~061 \\
P16 & 026, 064, 181 & $H_0$ is emergent (decompactified), not fundamental & Doc.~263, 309 \\textbf{[B]} \\
P17 & 028 & DE/DM relation circular; MOND partially derived & Doc.~028 (addendum) \\
P18 & 267 & Both pictures empirically indistinguishable; decision is theoretical & Doc.~267 \\
P20 & 267, 268 & Form: dimensionless $\xi$-ratio ($R_H/\ell_P$) is fixed & Docs.~267, 268 \\
P29 & 268 & Sequence $\{1,6,14,26\}$ read from data, not derived forward & Doc.~268, open \textbf{[S]} \\
P30 & 268 & Naive spherical Bessel sum does not reproduce $\{3,18,42,78\}$ & Doc.~268, open \textbf{[S]} \\
P31 & 268 & $T^4$ gives harmonic backbone $1:2:3:4:5$ (symmetric resonance $k^2=3n^2$) & Doc.~268 \textbf{[B]} \\
P32 & corpus-wide & $H_0$ is an imported external calibration, not derived & Doc.~309 \\
P33 & 009, 042, 181 & Form of $\xi$ forward: harmonic $4/3$, spatial dimensions $3$, $2^2=4$ from geometry & Doc.~009 \\
P34 & 181, 013 & 3+1 is empirical input; earlier ``follows necessarily'' claim overstated & Doc.~181 \\
P35 & corpus-wide & $\xi^N$ triviality: ``$X=\xi^N$'' is meaningless on its own & Docs.~182, 250 (addenda) \\
P36 & 182, cosm.\ sector & $R_H = c/H_0$ is the Hubble length, not the universe size & Doc.~182 \\
P37 & 009, 188, 271, 272, 275 & Three levels clearly separated: micro/meso/macro & Doc.~275 \textbf{[K]} \\
P38 & 275 & $1/\varphi$ is a transit point, not a fixed point & Doc.~275 \\
P39 & 182, 267, 277, corpus-wide & $H_0/R_H$ status: no framework derives from first principles & Doc.~309 \\
P40 & 005, 006, 011, 012, 133, 275 & Only ratios exact; absolute values carry fractal/SI correction via bridge constants & Doc.~A080 \textbf{[B]} \\
P41 & 267, 280, 221 & Checked, not booked (degeneracy risk) & Doc.~267 \\
P42 & 282, 292 & One declared reference point: $m_\mu/m_e$ & Doc.~292 \\
P43 & 006, 011, 292 & $K_{\text{frak}}$ not explicit in \texttt{calc\_En.py} & script corrected \\
P44 & 253, 190 & $\xi$ as linewidth: disproved for this Shor setup; window width is a measurement-setup property & Doc.~253 \\
R41 & 284 (cf.\ 270, 249) & Type-I decompactification is mode-preserving embedding (lossless) & Doc.~270 \textbf{[B]} \\
R42 & 270 (cf.\ 248) & Time is relational; symmetry breaking is the measurement process & Doc.~270 \textbf{[B]} \\
R43 & 285, 283 & $C_3<A_5$, $T^7=T^4\times T^3$: HLV $\supset$ FFGFT at group/dimension level & Doc.~285 \textbf{[B]} \\
R44 & 272, 285 & $\varphi$ acquires meaning as eigenvalue in HLV 5-fold/perp sector, not in FFGFT & Doc.~272 \textbf{[B]} \\
R45 & 285, 282--284 & Quasicrystal: structure/diffraction spectrum point-like; dynamical spectrum continuous & Doc.~285 \textbf{[B]} \\
R46 & 296 (cf.\ 025, 026, 063, 156) & Open test point: intermediate-scale formulation cross-cuts corpus line & open \textbf{[S]} \\
R47 & 285, 297, 298 & Spiral-time window not additive but internal (bounded excerpt) & Doc.~285 \textbf{[B]} \\
R48 & 304 (cf.\ 295, 301, 303) & Overstatement was in the mail only; document was correct & Doc.~304 \\
R49 & 304, 305 & Consolidated first-mention footnote for U1/U2/U3 and $M_\tau/M_B$ (Krüger 2026) & Docs.~304, 305 \\
R50 & 025, 063, 061, 064, 078 & Justification replaced; conclusion (no singular beginning; finite core $L_0$) unchanged & Doc.~306 \\
R51 & 049 (cf.\ 095, 306, 267, 302) & $T_0$ never denoted time zero, but the smallest possible time interval; notational legacy & Doc.~306 \\textbf{[B]} \\
R52 & 101, 165 & $L_0 = \xi\,\ell_P = 2.155\times10^{-39}$\,m; earlier values corrected & Doc.~180 \textbf{[K]} \\
R53 & 267 \S518 & Table row ``source of scale'': correctly ``set / calibrated to $H_0$'' & Doc.~267 \textbf{[B]} \\
R54 & 292, 306 & $N_0 \approx 38.6$ is a ladder-internal factor, not a fundamental unit; value open & Doc.~A105 \textbf{[K]} \\
R55 & 292 & Multiplicative structure: third digit arithmetically forced; evidence width one digit & Doc.~292 \textbf{[K]} \\
R56 & 188 \S & Three foundational assumptions declared: $\tilde{T}\cdot m=1$, $T^4$, $Z_3$ & Doc.~188 \textbf{[K]} \\
R57 & A015 & $\xi$ is a fundamental parameter; proof neither possible nor needed & Doc.~A015 \\
R58 & A142, A140 & Quantum gravity not required; $G=\xi^2/(4m_\text{char})$ intrinsic & Doc.~A142 \\
R59 & A130 & SI values scheme-dependent (1-loop); fractal corrections not separable & Doc.~A130 \textbf{[S]} \\
R60 & A160, A165 & $2\pi$ as circle circumference; $\xi/(2\pi)$ as phase correction per measurement & Docs.~A160, A165 \textbf{[B]} \\
R61 & 174, 175 & $\xi_\text{Higgs}$ selection data-driven; two-spaces doctrine withdrawn & Docs.~174, 175 (addenda) \\
R62 & 005 \S Baryon example & Anchor $\pi^2/2$, not $m_\mu$; error factor $46.705$ identified to $0.07\,\%$ & Doc.~A155 (candidate, open \\textbf{[S]}) \\
R63 & 253 \S $\xi$-resonance heuristic & $\xi$ inert in optimal range; optimisation disabled the filter & Doc.~253 \\
R64 & 2/python/t0\_no\_fallback\_shore.py & Procedure correct; separation plane too coarse; label wrong & Script corrected \\
R65 & 024, 075, 076, 147, 176, 253 & No procedure distinguishable from Shor; method claim withdrawn & Doc.~075 \textbf{[B]} \\
R66 & 097 & $16\pi^3$ counting correct, but not for the reason given in Doc.~097 & Doc.~310 \\
R67 & A040, A270 & Rounding interval leaves $n\in[98.3,\,102.0]$ open; $n=100$ not exactly forced & Doc.~A040 \\
R68 & 279, 308 & $\Lambda^{*}:=1/R_H^2$ as carrier of the $\Lambda$ function; name discipline like $a_0$ & Doc.~308 \\
R69 & 308 \S$a_0$ derivation & Compact time cycle $\tau_c=2\pi/H_0=91.9$\,Gyr; energy quantum $\hbar H_0$ & Doc.~308 \\
R70 & 061, A085 & $T_0$ scale gap and notation collision & Docs.~314, 315 \\
R71 & 309 & Sharpened: side-choice in the Hubble tension & Doc.~309 (addendum) \\
R72 & A010, A080, A265 & $E_0$: bare and corrected & Docs.~314, 315 \\
R73 & 026, 279 & $H_0$ notation, exponent $41/4$ & Docs.~026, 279 \\
R74 & 025, 063 & Lithium problem: mechanism rather than solution & Docs.~025, 063 \\
R75 & 024, 034, 075, 076, 147, 176 & Name collision $\xi$; new edition of Shor documents & Doc.~075 \textbf{[B]} \\
R76 & 318 & Validity range of mass derivations & Doc.~318 \\
R77 & 041, 160, 318 & Classification of early coupling-constant formulas & Doc.~041 (addendum) \\
R78 & 160, 005 & $\alpha_s$ and running coupling & Doc.~160 \\
R79 & 321, 319, 049 & Algebraic derivation of the $SU(3)_c$ gauge group & Docs.~319, 321 \\
R81 & 323, 321, 160, 320 & Weinberg angle at $M_Z$ from $\xi$ & Doc.~323 \\
R82 & 322, 320 & Hilbert space embedding and spectral theory & Doc.~322 \\
R83 & 319, 321, 323, 322, 320 & Status update of closed bridges & respective docs~\textbf{[B]} \\
R84 & 324, 321, 009 & Casimir derivation of $\xi$; $\xi$ not a G(6) spectral value & Doc.~324 \textbf{[K]} \\
R85 & 313, 325, 257, 321, 322 & Hawking mechanism microscopically closed & Doc.~325 \textbf{[B]} \\
R86 & 180, 322, 314, 321, 325 & Self-adjointness $\hat{F}$ complete; $\xi=\lambda_{\min}(\hat{F}_{D_4})$ & Doc.~327 \textbf{[B]} \\
R87 & 320, 322 & $\Delta m_{32}^2$: numerical values corrected ($+16.0\,\%$); closed by Doc.~340 ($1.0\,\%$) & Docs.~320, 340 \textbf{[K]} \\
R88 & 326, 329 & $n_{\mathrm{thresh}}$ scale-dependent; Matzke scale $=2\pi\ell_P/\ln 2$; universal number negative & Doc.~329 \textbf{[B]} \\
R89 & 180, 019, 201, 032 & $m_T = M_\mathrm{basis}/\xi$: one structure, sector-dependent basis mass & Doc.~180 \textbf{[B]} \\
R90 & 180, 019, 201 & Exponent in $L_0=\xi^1\ell_P$ derived from mass term, not by analogy & Doc.~180 \textbf{[B]} \\
R91 & 019, 201 & $\lambda$ in $m_T=\lambda/\xi$ is a basis mass, not the Higgs quartic & Docs.~019, 201 (addenda) \\
R92 & 250, 190-old & Doc.~250: $\xi/\ell$ is coupling $g_T$, not $m_T$; Schwarzschild probe for $L_0$ is an identity & Doc.~250; this new edition \\
R93 & 325, 329 & $S_\mathrm{BH}(M_\mathrm{coll})=2\pi$\,nat $=n_\mathrm{thresh}$(Matzke) in bits; structural constant & Docs.~325, 329 \textbf{[B]} \\
R94 & 325 & $I_\mathrm{sector}/I_\mathrm{thermal}=\log_2 3$ for all $M$ & Doc.~325 \textbf{[B]} \\
R95 & 330 & Two formulation clarifications; Type-I theorem [sketch]$\to$\textbf{[B]} after Krüger review & Doc.~330 \textbf{[B]} \\
R96 & 330, 332, 270 & Pullback: $L^2(S^1_m)\cong\mathrm{Per}_{R_m}(\mathbb{R}_t)$, not $L^2(\mathbb{R}_t)$ & Doc.~332 \textbf{[B]} \\
R97 & 324, 329 & Vacuum operator correction (Matzke): $V$ is involution in $\mathbb{Z}_3\mathbb{C}$ & Doc.~324 \textbf{[B]} \\
R98 & 333, 332, 315, 295, 285 & Two meanings of ``unrolling'' clarified; recursion tied to $\tilde{T}\cdot m$ & Doc.~333 \textbf{[B]} \\
R99 & 334, 307, 306, 285, 174 & Superposition without time: classical physics and QM inherit tacit instantaneity & Doc.~334 \textbf{[K]} \\
R100 & 336, 324, 329 & Notation shifted to symmetric $\mathbb{Z}_3\mathbb{C}$; Frobenius fixed points = FFGFT sectors & Doc.~336 \textbf{[B]} \\
R101 & 317, 336, 338 & $\xi=4/30000$ as Galois number; core identity $43200$ from GF$(9)^*$, GF$(27)$ & Docs.~317, 338 \textbf{[K]} \\
R102 & 338 & $1/\alpha=3700/27=137.037$ (7.6\,ppm) purely from Galois group orders & Doc.~338 \textbf{[K]} \\
R103 & 339, 340 & Frobenius decomposition GF$(27)^*$; neutrino mass hierarchy & Docs.~339, 340 \textbf{[K/B]} \\
R104 & 006, 338, 319, 340 & $v=248.3$\,GeV ($0.85\,\%$); $G_F$, $\tau_n$, $\Delta m$ derived & Doc.~319 \textbf{[K]} \\
R105 & 006, 011, 070, 182, 231, 257, 285, 306, 307 & Retroactive marker certification (before marker system, 1~Sept.~2026) & Doc.~006 (addendum) \\
R106 & 342, 343 & Prime forcing by GF$(3^k)$; physical primes $\{5,7,11,13\}$ forced & Doc.~342 \textbf{[B]} \\
R107 & 343, 338, 162, 186, 328 & $\xi$ as resolution floor; Farey coupling; resonance factoriser & Doc.~343 \textbf{[B/S]} \\
R108 & 340, 342, 343 & Dirac/Majorana decomposition of GF$(27)^*$ classes; $\nu_3$ = Dirac neutrino & Docs.~343, 340 \textbf{[B]} \\
R109 & 006 & Neutrino part of Doc.~006 not viable \textbf{[X]}; Yukawa method valid for charged leptons & Doc.~006 (addendum) \\
R110 & 320, 022 & Addendum boxes: $\nu_3$ Dirac neutrino; $\xi^3$ suppression not a sterile signal & Docs.~320, 022 \\
R111 & 344--348 & Charge quantisation, Gell-Mann-Nishijima and generations from GF$(27)^*$ & Docs.~344--348 \textbf{[B]} \\
R112 & 006, 319, 351 & $v=246$\,GeV is SM definition, not a measurement \QM & Doc.~351 \textbf{[K]} \\
R113 & 338, 317, 011, 351 & Anchor $m_\mu/m_e$ is QED-dependent; comparison values not theory-free \QM & Doc.~351 \textbf{[K]} \\
R114 & 351, 352 & Lepton residuals: Koide most precise ($0.046\cdot\xi$); Galois residual $-77.6\cdot\xi$; Koide amplitude $d/c=\sqrt{2}$ closed by Doc.~353 \textbf{[B]} & Doc.~352 \textbf{[K]}; Doc.~353 \textbf{[B]} \\
K-041-a & 041 \S$\delta_\mathrm{CKM}$ & Formula $\arcsin(2\sqrt{2}\cdot\sqrt{\xi}/3)$ gives $0.011$\,rad; table value $1.20$\,rad; candidate: $\arcsin(2\sqrt{2}/3)$ & Doc.~041 (addendum) \\
K-041-b & 041 \S$\theta_{13}$ & Formula $\arcsin(\xi^{1/3})$ gives $2.93°$; table value $8.57°$; candidate: $1/300=25\xi$ & Doc.~041 (addendum) \\
K-041-c & 041 \S$|V_{us}|$ & Formula gives $0.2124$; table value $0.22452$ ($\Delta=5.4\,\%$) & Doc.~041 (addendum) \\
R115 & 355, 348 & GF$(3^6)$=GF$(729)$: 7th roots of unity require degree-6 extension; $728=8\times7\times13$; candidate for CKM/PMNS angles & Doc.~355 (note); open \textbf{[S]} \\
R116 & 356, 357 & SM overview sheets: Doc.~356 particle diagram (all fermion quantum numbers from GF$(27)^*$, landscape A4); Doc.~357 certificate table of all quantum numbers with status; Koide amplitude $d/c=\sqrt{2}$ entered as \textbf{[B]} (7~Sept.~2026) & Docs.~356, 357 \textbf{[B]} \\
R117 & 358 & Harmonics and algebraic structure --- why the same mathematics; Thm~A \textbf{[B]}: no comma in finite field; Thm~B \textbf{[B]}: all nine Yukawa coefficients in Galois grid $\{2,3,5,7,11,13\}$; Obs~C \textbf{[S]}: three twelves; Thm~D \textbf{[B]}: harmonics $=$ two representations of $T^4/Z_3$; script 22/22 & Doc.~358 \textbf{[B/S]} \\\
R118 & 359 & AI-based pattern recognition and algebraic limits: HRV band limits in 5-limit grid \textbf{[B]}; resolution limit $100\xi$ appears empirically in AI models \textbf{[B/L]}; G-CNNs for GF$(3^k)$ missing (gap identified); PAC non-learnability of admissibility criterion \textbf{[B]}; Galois-informed network sketch \textbf{[S]}; script 14/14 & Doc.~359 \textbf{[B/S]} \\\
R119 & 360 & Galois-HRV preprocessing and orbit experiment; Galois stability separates severe VES from normal significantly (Mann-Whitney $p<0.001$); orbit-$k$ alone not discriminative; normal range LF/HF$<4$ fits Galois grid $k=1..4$; next free Doc. No. 361 & Doc.~360 	extbf{[S/K]} \
\end{longtable}
\endgroup






